\subsection{Plasma defocusing}
Several effects are at the origin of electrons in the conduction band : thermal excitation, tunneling ionization, multiphoton ionization, free electrons resulting from cristal defects. All of these free carriers reacts to the incident electric field. The simplest modelisation is to treat the conduction electrons as a free-electron gas subject to the friction
 $\gamma = 1/\tau$ of Eq.~\eqref{eq:tau_result}, the equation of motion in a field
 $\vec{E}(t) = \vec{E}_0\, e^{-i\omega_0 t}$ reads
\begin{equation}
  m^\ast \frac{d^2}{dt^2}\vec{x} = -e\vec{E} - \frac{m^\ast}{\tau}\, \frac{d}{dt}\vec{x}
  \,\Longrightarrow\,
  \vec{x}
  = \frac{\tau e\vec{E}}{m^\ast \omega_0\left(\omega_0\tau + i\right)}
\end{equation}
The polarisation of the free carriers $\vec{P}_{\mathrm{free}} = -eN\vec{x}$ follows as
\begin{equation}
  \vec{P}_{\mathrm{free}} = -eN\vec{x}
 =-\frac{\omega_p^2\tau}{ \omega_0\left(\omega_0\tau + i\right)}\varepsilon_0\vec{E}
\end{equation}
Introducing the plasma frequency
 $\omega_p^2 = \dfrac{Ne^2}{\varepsilon_0 m^\ast}$.

The valence electrons are not removed by the excitation, and their polarisation
 $\vec{P}_{\mathrm{bound}} = \varepsilon_0 \chi^{(1)} \vec{E}$ must be retained. Both
populations are driven by the same macroscopic field and respond linearly, so their
contributions add.Introducing the linear index
 $n_0^2 = 1 + \chi^{(1)}$ of the unexcited medium, so that the relative permitivity read : 
\begin{equation}
  \varepsilon_r = \frac{D}{\varepsilon_0 E} = \frac{1}{\varepsilon_0\vec{E}}\left(\varepsilon_0 \vec{E} + \vec{P}_{\mathrm{bound}} + \vec{P}_{\mathrm{free}}\right)
     =  n_0^2
    - \frac{\omega_p^2\tau}{\omega_0(\omega_0\tau + i)}
\end{equation}
multiplying numerator and denominator by the complex conjugate and separates the real and imaginary parts we obtain 
\begin{align}
\varepsilon_r &= \varepsilon_1 + i\varepsilon_2\\
  \varepsilon_1 &= n_0^2 - \frac{\omega_p^2\tau^2}{1 + \omega_0^2\tau^2} , \\
  \varepsilon_2 &= \frac{\omega_p^2\tau}{\omega_0\left(1 + \omega_0^2\tau^2\right)} > 0 ,
\end{align}

The complex index follows from $\tilde{n} = n + i\kappa = \sqrt{\varepsilon_r}$, so that
 $\varepsilon_1 = n^2 - \kappa^2$ and $\varepsilon_2 = 2n\kappa$.

In the case here the number of oscillations before collision is high, the real part of the
refractive index simplify
\begin{equation}
  n = \sqrt{\varepsilon_1} = \left(n_0^2 - \frac{\omega_p^2\tau^2}{1 + \omega_0^2\tau^2}\right)^\frac{1}{2} \approx n_0\left(1 - \frac{\omega_p^2}{n_0^2\,\omega_0^2}\right)^\frac{1}{2}
\end{equation}
The square root must be real, which imposes
 $\dfrac{\omega_p^2}{n_0^2\,\omega_0^2} \leq 1$. Writing out the definition of the plasma
frequency,
\begin{equation}
  \frac{Ne^2}{n_0^2\,\varepsilon_0 m^\ast \omega_0^2} \leq 1
  \,\Longrightarrow\,
  N \leq \frac{n_0^2\,\varepsilon_0 m^\ast \omega_0^2}{e^2}
\end{equation}
Defining the critical density $N_c$ as the density at which the plasma frequency equals the
laser frequency, $\omega_p(N_c) = \omega_0$,
\begin{equation}
  N_c = \frac{\varepsilon_0 m^\ast \omega_0^2}{e^2}
  \,\Longrightarrow\,
  \frac{\omega_p^2}{\omega_0^2} = \frac{N}{N_c}
\end{equation}
the condition becomes $N \leq n_0^2 N_c$, beyond which the wave is evanescent and the
plasma is reflective. The index is then
\begin{equation}
  n = n_0\sqrt{1 - \frac{\omega_p^2}{n_0^2\,\omega_0^2}}
    = n_0\sqrt{1 - \frac{N}{n_0^2 N_c}}
\end{equation}
In the case of an underdense plasma, $N \ll n_0^2 N_c$, the real part of the refractive
index can be expanded as
\begin{equation}
  n \simeq n_0 - \frac{N}{2\,n_0 N_c}
  \label{eq:delta_n_plasma}
\end{equation}
